
\documentclass{article}
\usepackage{colortbl}
\usepackage{makecell}
\usepackage{multirow}
\usepackage{supertabular}

\begin{document}

\newcounter{utterance}

\centering \large Interaction Transcript for game `hot\_air\_balloon', experiment `air\_balloon\_survival\_fi\_complexity\_hard', episode 1 with mistral{-}large{-}3{-}azure.
\vspace{24pt}

{ \footnotesize  \setcounter{utterance}{1}
\setlength{\tabcolsep}{0pt}
\begin{supertabular}{c@{$\;$}|p{.15\linewidth}@{}p{.15\linewidth}p{.15\linewidth}p{.15\linewidth}p{.15\linewidth}p{.15\linewidth}}
    \# & $\;$A & \multicolumn{4}{c}{Game Master} & $\;\:$B\\
    \hline

    \theutterance \stepcounter{utterance}  
    & & \multicolumn{4}{p{0.6\linewidth}}{
        \cellcolor[rgb]{0.9,0.9,0.9}{
            \makecell[{{p{\linewidth}}}]{
                \texttt{\tiny{[P1$\langle$GM]}}
                \texttt{Osallistut yhteistoiminnalliseen neuvottelupeliin.} \\
\\ 
\texttt{Teidän on yhdessä toisen osallistujan kanssa sovittava yhdestä joukosta säilytettäviä esineitä. Kummallakin teistä on oma näkemyksensä siitä, kuinka paljon kukin esine on teille tärkeä (tärkeys). Ette tiedä, miten toinen osallistuja arvostaa esineitä. Lisäksi teille ilmoitetaan, kuinka paljon vaivaa kukin esine vaatii.} \\
\texttt{Voitte sopia joukosta vain, jos valittujen esineiden yhteenlaskettu vaiva ei ylitä yhteistä rajaa:} \\
\\ 
\texttt{LIMIT = 4497} \\
\\ 
\texttt{Tässä ovat yksittäisten esineiden vaiva{-}arvot:} \\
\\ 
\texttt{Esineen vaiva = \{"C08": 828, "C55": 250, "A43": 412, "A67": 75, "B42": 275, "A62": 565, "B19": 889, "C43": 73, "C73": 747, "C10": 77, "B36": 23, "A18": 651, "C61": 11, "C67": 298, "C42": 769, "C70": 812, "B40": 368, "C39": 506, "A21": 481, "C16": 884\}} \\
\\ 
\texttt{Tässä on henkilökohtainen näkemyksesi kunkin esineen tärkeydestä:} \\
\\ 
\texttt{Esineen tärkeysarvot = \{"C08": 916, "C55": 338, "A43": 500, "A67": 163, "B42": 363, "A62": 653, "B19": 977, "C43": 161, "C73": 835, "C10": 165, "B36": 111, "A18": 739, "C61": 99, "C67": 386, "C42": 857, "C70": 900, "B40": 456, "C39": 594, "A21": 569, "C16": 972\}} \\
\\ 
\texttt{Tavoite:} \\
\\ 
\texttt{Tavoitteenasi on neuvotella yhteinen joukko esineitä, jotka hyödyttävät sinua mahdollisimman paljon (eli maksimoivat kokonaistärkeyden SINULLE) ja pysyvät samalla RAJAN sisällä. Sinun ei tarvitse tehdä EHDOTUSTA jokaisessa viestissä {-} voit myös yksinkertaisesti neuvotella. Kaikki taktiikat ovat sallittuja!} \\
\\ 
\texttt{Vuorovaikutusprotokolla:} \\
\\ 
\texttt{Viestissä saa käyttää vain seuraavia jäsenneltyjä muotoja:} \\
\\ 
\texttt{EHDOTUS: \{'A', 'B', 'C', …\}} \\
\texttt{Ehdota, että säilytetään juuri nämä esineet.} \\
\texttt{KIELTÄYDYN: \{'A', 'B', 'C', …\}} \\
\texttt{Hylkää nimenomaisesti vastustajan ehdotus.} \\
\texttt{ARGUMENTTI: \{'...'\}} \\
\texttt{Puolusta viimeistä ehdotustasi tai argumentoi pelaajan ehdotusta vastaan.} \\
\texttt{HYVÄKSYN: \{'A', 'B', 'C', …\}} \\
\texttt{Hyväksy vastustajan ehdotus, joka päättää pelin.} \\
\\ 
\\ 
\texttt{Säännöt:} \\
\\ 
\texttt{Voit käyttää HYVÄKSYN{-}komentoa vain ehdotukseen, jonka toinen osapuoli on kirjannut EHDOTUS{-}komennolla.} \\
\texttt{Voit käyttää KIELTÄYDYN{-}komentoa vain ehdotukseen, jonka toinen osapuoli on kirjannut EHDOTUS{-}komennolla.} \\
\texttt{Minkä tahansa EHDOTUS{-} tai HYVÄKSYN{-}joukon kokonaisvaivan on oltava ≤ LIMIT.} \\
\texttt{ÄLÄ paljasta piilotettuja tärkeysarvojasi.} \\
\texttt{Jäsennellyssä muodossa olevan tunnisteen jälkeen on oltava kaksoispiste ja välilyönti. Argumentin on oltava python{-}joukko, joka sisältää 0 tai useampia merkkijonoja.} \\
\texttt{Sen on siis oltava muotoa "TAG: \{...\}"} \\
\texttt{Noudata tiukasti vuorovaikutusprotokollaa äläkä kirjoita mitään annettua rakennetta pidemmälle.} \\
\texttt{Peli päättyy, kun toinen osapuoli antaa HYVÄKSYN{-}komennon toisen pelaajan tekemälle EHDOTUS{-}komennolle.} \\
\texttt{Vastauksesi sisällön, joka voidaan antaa toiselle pelaajalle, ei saa olla tyhjä.} \\
\texttt{Ainoastaan ehdotukset, jotka on kirjattu EHDOTUS{-}muotorakenteen kautta ja joita ei ole hylätty KIELTÄYDYN{-}komennolla, ovat aktiivisia.} \\
\\ 
            }
        }
    }
    & \\ \\

    \theutterance \stepcounter{utterance}  
    & \multicolumn{4}{p{0.6\linewidth}}{
        \cellcolor[rgb]{0.9,0.9,0.9}{
            \makecell[{{p{\linewidth}}}]{
                \texttt{\tiny{[P1$\rangle$GM]}}
                \texttt{EHDOTUS: \{'C08', 'C55', 'A43', 'A67', 'B42', 'A62', 'C43', 'C73', 'C10', 'B36', 'C61', 'C67', 'B40', 'C39', 'A21'\}} \\
            }
        }
    }
    & & \\ \\

    \theutterance \stepcounter{utterance}  
    & & \multicolumn{4}{p{0.6\linewidth}}{
        \cellcolor[rgb]{0.9,0.9,0.9}{
            \makecell[{{p{\linewidth}}}]{
                \texttt{\tiny{[GM$\rangle$P2]}}
                \texttt{Osallistut yhteistoiminnalliseen neuvottelupeliin.} \\
\\ 
\texttt{Teidän on yhdessä toisen osallistujan kanssa sovittava yhdestä joukosta säilytettäviä esineitä. Kummallakin teistä on oma näkemyksensä siitä, kuinka paljon kukin esine on teille tärkeä (tärkeys). Ette tiedä, miten toinen osallistuja arvostaa esineitä. Lisäksi teille ilmoitetaan, kuinka paljon vaivaa kukin esine vaatii.} \\
\texttt{Voitte sopia joukosta vain, jos valittujen esineiden yhteenlaskettu vaiva ei ylitä yhteistä rajaa:} \\
\\ 
\texttt{LIMIT = 4497} \\
\\ 
\texttt{Tässä ovat yksittäisten esineiden vaivat:} \\
\\ 
\texttt{Esineen vaiva = \{"C08": 828, "C55": 250, "A43": 412, "A67": 75, "B42": 275, "A62": 565, "B19": 889, "C43": 73, "C73": 747, "C10": 77, "B36": 23, "A18": 651, "C61": 11, "C67": 298, "C42": 769, "C70": 812, "B40": 368, "C39": 506, "A21": 481, "C16": 884\}} \\
\\ 
\texttt{Tässä on henkilökohtainen näkemyksesi kunkin esineen tärkeydestä:} \\
\\ 
\texttt{Esineen tärkeysarvot = \{"C08": 916, "C55": 338, "A43": 500, "A67": 163, "B42": 363, "A62": 653, "B19": 977, "C43": 161, "C73": 835, "C10": 165, "B36": 111, "A18": 739, "C61": 99, "C67": 386, "C42": 857, "C70": 900, "B40": 456, "C39": 594, "A21": 569, "C16": 972\}} \\
\\ 
\texttt{Tavoite:} \\
\\ 
\texttt{Tavoitteenasi on neuvotella yhteinen joukko esineitä, jotka hyödyttävät sinua mahdollisimman paljon (eli maksimoivat kokonaistärkeyden SINULLE) ja pysyvät samalla RAJAN sisällä. Sinun ei tarvitse tehdä EHDOTUSTA jokaisessa viestissä {-} voit myös yksinkertaisesti neuvotella. Kaikki taktiikat ovat sallittuja!} \\
\\ 
\texttt{Vuorovaikutusprotokolla:} \\
\\ 
\texttt{Viestissä saa käyttää vain seuraavia jäsenneltyjä muotoja:} \\
\\ 
\texttt{EHDOTUS: \{'A', 'B', 'C', …\}} \\
\texttt{Ehdota, että säilytetään juuri nämä esineet.} \\
\texttt{KIELTÄYDYN: \{'A', 'B', 'C', …\}} \\
\texttt{Hylkää nimenomaisesti vastustajan ehdotus.} \\
\texttt{ARGUMENTTI: \{'...'\}} \\
\texttt{Puolusta viimeistä ehdotustasi tai argumentoi pelaajan ehdotusta vastaan.} \\
\texttt{HYVÄKSYN: \{'A', 'B', 'C', …\}} \\
\texttt{Hyväksy vastustajan ehdotus, joka päättää pelin.} \\
\\ 
\\ 
\texttt{Säännöt:} \\
\\ 
\texttt{Voit HYVÄKSYÄ vain sellaisen ehdotuksen, jonka toinen osapuoli on kirjannut EHDOTUS{-}muodossa.} \\
\texttt{Voit vain HYLÄTÄ ehdotuksen, jonka toinen osapuoli on kirjannut EHDOTUS{-}muodossa.} \\
\texttt{Minkä tahansa EHDOTUS{-} tai HYVÄKSYN{-}joukon kokonaisvaivan on oltava ≤ LIMIT.} \\
\texttt{ÄLÄ paljasta piilotettuja tärkeysarvojasi.} \\
\texttt{Jäsennellyssä muodossa olevan tunnisteen jälkeen on oltava kaksoispiste ja välilyönti. Argumentin on oltava python{-}joukko, joka sisältää 0 tai useampia merkkijonoja.} \\
\texttt{Sen on siis oltava muotoa "TAG: \{...\}"} \\
\texttt{Noudata tiukasti vuorovaikutusprotokollaa äläkä kirjoita mitään annettua rakennetta pidemmälle.} \\
\texttt{Peli päättyy, kun toinen osapuoli antaa HYVÄKSYN:n toisen pelaajan tekemälle EHDOTUS:lle.} \\
\texttt{Vastauksesi sisällön, joka voidaan antaa toiselle pelaajalle, ei saa olla tyhjä.} \\
\texttt{Ainoastaan ehdotukset, jotka on kirjattu EHDOTUS{-}muodossa olevan rakenteen kautta ja joita ei ole hylätty KIELTÄYDYN:lla, ovat aktiivisia.} \\
\\ 
\\ 
\\ 
\texttt{Saat nyt toisen pelaajan ensimmäisen viestin.} \\
\\ 
\texttt{EHDOTUS: \{'C08', 'C55', 'A43', 'A67', 'B42', 'A62', 'C43', 'C73', 'C10', 'B36', 'C61', 'C67', 'B40', 'C39', 'A21'\}} \\
            }
        }
    }
    & \\ \\

    \theutterance \stepcounter{utterance}  
    & & & \multicolumn{4}{p{0.6\linewidth}}{
        \cellcolor[rgb]{0.9,0.9,0.9}{
            \makecell[{{p{\linewidth}}}]{
                \texttt{\tiny{[GM$\langle$P2]}}
                \texttt{KIELTÄYDYN: \{'C08', 'C55', 'A43', 'A67', 'B42', 'A62', 'C43', 'C73', 'C10', 'B36', 'C61', 'C67', 'B40', 'C39', 'A21'\}} \\
\\ 
\texttt{ARGUMENTTI: \{'Tämä ehdotus ylittää vaivalimittimme, sillä sen kokonaisvaiva on 5882, mikä on yli LIMIT{-}arvon 4497. Lisäksi ehdotuksessa on useita esineitä, jotka eivät ole minulle erityisen tärkeitä (esim. C43, C10, B36). Ehdotan, että keskitymme ensisijaisesti esineisiin, joilla on korkea tärkeys minulle ja joiden vaiva on kohtuullinen.'\}} \\
            }
        }
    }
     \\ \\

    \theutterance \stepcounter{utterance}  
    & & \multicolumn{4}{p{0.6\linewidth}}{
        \cellcolor[rgb]{0.9,0.9,0.9}{
            \makecell[{{p{\linewidth}}}]{
                \texttt{\tiny{[P1$\langle$GM]}}
                \texttt{KIELTÄYDYN: \{'C08', 'C55', 'A43', 'A67', 'B42', 'A62', 'C43', 'C73', 'C10', 'B36', 'C61', 'C67', 'B40', 'C39', 'A21'\}} \\
\\ 
\texttt{ARGUMENTTI: \{'Tämä ehdotus ylittää vaivalimittimme, sillä sen kokonaisvaiva on 5882, mikä on yli LIMIT{-}arvon 4497. Lisäksi ehdotuksessa on useita esineitä, jotka eivät ole minulle erityisen tärkeitä (esim. C43, C10, B36). Ehdotan, että keskitymme ensisijaisesti esineisiin, joilla on korkea tärkeys minulle ja joiden vaiva on kohtuullinen.'\}} \\
            }
        }
    }
    & \\ \\

    \theutterance \stepcounter{utterance}  
    & \multicolumn{4}{p{0.6\linewidth}}{
        \cellcolor[rgb]{0.9,0.9,0.9}{
            \makecell[{{p{\linewidth}}}]{
                \texttt{\tiny{[P1$\rangle$GM]}}
                \texttt{EHDOTUS: \{'C08', 'A62', 'B19', 'C73', 'C42', 'C70', 'A18', 'A21', 'C16'\}} \\
            }
        }
    }
    & & \\ \\

    \theutterance \stepcounter{utterance}  
    & & \multicolumn{4}{p{0.6\linewidth}}{
        \cellcolor[rgb]{0.9,0.9,0.9}{
            \makecell[{{p{\linewidth}}}]{
                \texttt{\tiny{[GM$\rangle$P2]}}
                \texttt{EHDOTUS: \{'C08', 'A62', 'B19', 'C73', 'C42', 'C70', 'A18', 'A21', 'C16'\}} \\
            }
        }
    }
    & \\ \\

    \theutterance \stepcounter{utterance}  
    & & & \multicolumn{4}{p{0.6\linewidth}}{
        \cellcolor[rgb]{0.9,0.9,0.9}{
            \makecell[{{p{\linewidth}}}]{
                \texttt{\tiny{[GM$\langle$P2]}}
                \texttt{KIELTÄYDYN: \{'C08', 'A62', 'B19', 'C73', 'C42', 'C70', 'A18', 'A21', 'C16'\}} \\
\\ 
\texttt{ARGUMENTTI: \{'Tämä ehdotus on liian raskas: kokonaisvaiva on 6218, mikä ylittää selvästi LIMIT{-}arvon 4497. Lisäksi ehdotuksessa on useita korkean vaivan esineitä (esim. B19, C70, C16), jotka voivat syödä liikaa tilaa. Ehdotan, että priorisoimme esineitä, joiden vaiva{-}tärkeys{-}suhde on parempi. Voisimme esimerkiksi jättää pois B19:n (vaiva 889) ja C16:n (vaiva 884), jolloin saadaan enemmän tilaa muille tärkeille esineille.'\}} \\
\\ 
\texttt{**Uusi ehdotukseni:**} \\
\texttt{EHDOTUS: \{'C08', 'A62', 'C73', 'C42', 'C70', 'A18', 'A21'\} *(Kokonaisvaiva: 4412 ≤ LIMIT)*} \\
            }
        }
    }
     \\ \\

    \theutterance \stepcounter{utterance}  
    & & \multicolumn{4}{p{0.6\linewidth}}{
        \cellcolor[rgb]{0.9,0.9,0.9}{
            \makecell[{{p{\linewidth}}}]{
                \texttt{\tiny{[GM$\rangle$P2]}}
                \texttt{Käytit EHDOTUS{-}tunnistetta, mutta et antanut kelvollista python{-}joukkoa, joka sisältää merkkijonoja argumentteina, esim. \{'A', 'B', 'C', ...\}. Yritä uudelleen.} \\
            }
        }
    }
    & \\ \\

    \theutterance \stepcounter{utterance}  
    & & & \multicolumn{4}{p{0.6\linewidth}}{
        \cellcolor[rgb]{0.9,0.9,0.9}{
            \makecell[{{p{\linewidth}}}]{
                \texttt{\tiny{[GM$\langle$P2]}}
                \texttt{EHDOTUS: \{'C08', 'A62', 'C73', 'C42', 'C70', 'A18', 'A21'\}} \\
\\ 
\texttt{ARGUMENTTI: \{'Tämä ehdotus on LIMIT:in sisällä (vaiva 4412) ja sisältää minulle tärkeitä esineitä. Jos haluat lisätä muita esineitä, voimme harkita esimerkiksi B42:n (vaiva 275) tai C55:n (vaiva 250) poisjättämistä, jotta saadaan tilaa pienemmille esineille.'\}} \\
            }
        }
    }
     \\ \\

    \theutterance \stepcounter{utterance}  
    & & \multicolumn{4}{p{0.6\linewidth}}{
        \cellcolor[rgb]{0.9,0.9,0.9}{
            \makecell[{{p{\linewidth}}}]{
                \texttt{\tiny{[P1$\langle$GM]}}
                \texttt{EHDOTUS: \{'C08', 'A62', 'C73', 'C42', 'C70', 'A18', 'A21'\}} \\
\\ 
\texttt{ARGUMENTTI: \{'Tämä ehdotus on LIMIT:in sisällä (vaiva 4412) ja sisältää minulle tärkeitä esineitä. Jos haluat lisätä muita esineitä, voimme harkita esimerkiksi B42:n (vaiva 275) tai C55:n (vaiva 250) poisjättämistä, jotta saadaan tilaa pienemmille esineille.'\}} \\
            }
        }
    }
    & \\ \\

    \theutterance \stepcounter{utterance}  
    & \multicolumn{4}{p{0.6\linewidth}}{
        \cellcolor[rgb]{0.9,0.9,0.9}{
            \makecell[{{p{\linewidth}}}]{
                \texttt{\tiny{[P1$\rangle$GM]}}
                \texttt{KIELTÄYDYN: \{'C08', 'A62', 'C73', 'C42', 'C70', 'A18', 'A21'\}} \\
\\ 
\texttt{ARGUMENTTI: \{'Ehdotuksesi kokonaisvaiva (4412) on hyväksyttävällä tasolla, mutta haluaisin priorisoida esineitä, jotka ovat minulle erityisen arvokkaita. Voisimme esimerkiksi korvata C70:n (vaiva 812) ja A21:n (vaiva 481) esineillä B42 (vaiva 275) ja C55 (vaiva 250), jolloin saadaan tilaa lisäesineille. Tämä parantaisi minulle tärkeiden esineiden kattavuutta ilman, että ylitetään LIMIT.'\}} \\
\\ 
\texttt{EHDOTUS: \{'C08', 'A62', 'B19', 'C73', 'C42', 'A18', 'B42', 'C55'\}} \\
            }
        }
    }
    & & \\ \\

    \theutterance \stepcounter{utterance}  
    & & \multicolumn{4}{p{0.6\linewidth}}{
        \cellcolor[rgb]{0.9,0.9,0.9}{
            \makecell[{{p{\linewidth}}}]{
                \texttt{\tiny{[GM$\rangle$P2]}}
                \texttt{KIELTÄYDYN: \{'C08', 'A62', 'C73', 'C42', 'C70', 'A18', 'A21'\}} \\
\\ 
\texttt{ARGUMENTTI: \{'Ehdotuksesi kokonaisvaiva (4412) on hyväksyttävällä tasolla, mutta haluaisin priorisoida esineitä, jotka ovat minulle erityisen arvokkaita. Voisimme esimerkiksi korvata C70:n (vaiva 812) ja A21:n (vaiva 481) esineillä B42 (vaiva 275) ja C55 (vaiva 250), jolloin saadaan tilaa lisäesineille. Tämä parantaisi minulle tärkeiden esineiden kattavuutta ilman, että ylitetään LIMIT.'\}} \\
\\ 
\texttt{EHDOTUS: \{'C08', 'A62', 'B19', 'C73', 'C42', 'A18', 'B42', 'C55'\}} \\
            }
        }
    }
    & \\ \\

    \theutterance \stepcounter{utterance}  
    & & & \multicolumn{4}{p{0.6\linewidth}}{
        \cellcolor[rgb]{0.9,0.9,0.9}{
            \makecell[{{p{\linewidth}}}]{
                \texttt{\tiny{[GM$\langle$P2]}}
                \texttt{KIELTÄYDYN: \{'C08', 'A62', 'B19', 'C73', 'C42', 'A18', 'B42', 'C55'\}} \\
\\ 
\texttt{ARGUMENTTI: \{'B19:n (vaiva 889) sisällyttäminen jättää liian vähän tilaa muille tärkeille esineille. Sen poisjättämällä voisimme lisätä esimerkiksi C16:n (vaiva 884), joka on minulle erittäin tärkeä. Ehdotan seuraavaa joukkoa, jossa on korkeampi kokonaistärkeys minulle ja joka pysyy LIMIT:in sisällä:'\}} \\
\\ 
\texttt{EHDOTUS: \{'C08', 'A62', 'C73', 'C42', 'C16', 'A18', 'B42', 'C55'\} *(Kokonaisvaiva: 4414 ≤ LIMIT)*} \\
            }
        }
    }
     \\ \\

    \theutterance \stepcounter{utterance}  
    & & & \multicolumn{2}{p{0.3\linewidth}}{
        \cellcolor[rgb]{0.9,0.9,0.9}{
            \makecell[{{p{\linewidth}}}]{
                \texttt{\tiny{[GM$|$GM]}}
                \texttt{abort} \\
            }
        }
    }
    & & \\ \\

    \theutterance \stepcounter{utterance}  
    & & & \multicolumn{2}{p{0.3\linewidth}}{
        \cellcolor[rgb]{0.9,0.9,0.9}{
            \makecell[{{p{\linewidth}}}]{
                \texttt{\tiny{[GM$|$GM]}}
                \texttt{end game} \\
            }
        }
    }
    & & \\ \\

\end{supertabular}
}

\end{document}
